% =====================================================================
% Methodology
% =====================================================================
\section{Methodology}
\label{sec:methodology}

\subsection{Implementation overview}
\label{sec:impl-overview}

All four interpolation methods were implemented in Python~3.12 using
NumPy~\cite{numpy} and SciPy~\cite{scipy}. As the proposal required,
each method was coded directly from the textbook formulation rather
than wrapping a library black-box. SciPy was used only as a numerical
linear-algebra back-end---specifically \texttt{scipy.linalg.solve\_banded}
for the cubic-spline tridiagonal system and \texttt{scipy.linalg.solve}
for the quintic-spline dense system. The B-spline basis was built via
the Cox--de Boor recursion (Equation~\ref{eq:cox-de-boor}) without
relying on \texttt{scipy.interpolate.BSpline}. The complete source
code, datasets, and results are available in the project repository.

The cubic and B-spline implementations were validated against
SciPy's reference \texttt{CubicSpline} and \texttt{BSpline} classes
on identical inputs and agreed to machine precision (maximum error
$\sim 10^{-15}$). The quintic-spline implementation was validated by
reconstructing a known fifth-degree polynomial; the implementation
reproduced both the polynomial values and its analytic derivatives
to floating-point precision ($\sim 10^{-13}$).

\subsection{Datasets}
\label{sec:datasets}

\subsubsection{Mobile robot scenario.}
The mobile-robot dataset consists of twelve $(x, y)$ waypoints in
metres representing a robot navigating a warehouse aisle. The path
includes a straight entry segment, a sharp 90$^{\circ}$ right turn
around a shelving unit, a diagonal traverse to a pick station, a
90$^{\circ}$ left turn back into the main aisle, and a gentle
S-curve to the exit. Waypoint times use chord-length parameterisation
at a nominal speed of 1\,m/s, the standard recommendation in
\cite[Section 24.1]{chapra2021} for non-uniformly spaced data. Total
trajectory time is 24.93\,s.

\subsubsection{Robot arm scenario.}
The robot-arm dataset describes a 2-DOF planar manipulator with link
lengths $L_1 = \SI{0.5}{\meter}$ and $L_2 = \SI{0.4}{\meter}$
performing a pick-and-place cycle. The cycle visits eight joint-space
waypoints (in degrees, converted to radians internally): home,
approach over part, grasp, lift, transfer mid-point, over placement
bay, place, retract. Forward kinematics maps joint angles to
Cartesian end-effector position via
\begin{align}
x &= L_1 \cos\theta_1 + L_2 \cos(\theta_1 + \theta_2),\\
y &= L_1 \sin\theta_1 + L_2 \sin(\theta_1 + \theta_2),
\end{align}
following Chapter~3 of \cite{craig2018}. All Cartesian waypoints lie
within the reachable workspace $\|(x,y)\| \le L_1 + L_2 = \SI{0.9}{\meter}$.
Waypoint timing is uniform at one second per segment for a total of
seven seconds.

\subsubsection{Variant datasets for the experimental sweep.}
To investigate how each method's behaviour depends on input
characteristics, we generated five additional dataset variants by
sweeping over three independent axes: waypoint density, geometric
character, and time spacing
(\cref{tab:variant-summary}). The variant family is designed so that
each axis varies in isolation: the density sweep holds geometry and
spacing fixed, the geometry sweep holds density and spacing fixed,
and the spacing sweep holds geometry and density fixed.

\begin{table}[!t]
\centering
\caption{Dataset variants used in the experimental sweep. The
canonical baseline is \texttt{mixed\_baseline\_chord}; each row
varies exactly one axis from the baseline.}
\label{tab:variant-summary}
\begin{tabular}{lccc}
\toprule
\textbf{Variant} & \textbf{Geometry} & \textbf{Spacing} & \textbf{$n$} \\
\midrule
mixed\_sparse\_chord       & mixed   & chord   &  5 \\
mixed\_baseline\_chord     & mixed   & chord   & 12 \\
mixed\_dense\_chord        & mixed   & chord   & 20 \\
sharp\_n12\_chord          & sharp   & chord   & 12 \\
gradual\_n12\_chord        & gradual & chord   & 12 \\
mixed\_baseline\_uniform   & mixed   & uniform & 12 \\
robot\_arm\_baseline       & arm     & uniform &  8 \\
\bottomrule
\end{tabular}
\end{table}

The geometry options are: a \emph{mixed} warehouse-aisle path
(arc-length resampled from the original 12-vertex polyline at the
desired density); a \emph{sharp} rectangular zigzag with eleven
segments alternating between horizontal and vertical and exact 90$^{\circ}$
interior corners; and a \emph{gradual} sinusoidal path
$y = A \sin(2\pi k x / X)$ chosen so every method should agree
closely. The spacing options are chord-length parameterisation and
uniform time spacing.

\subsection{Evaluation metrics}
\label{sec:metrics}

We evaluate each method-dataset pairing using the metrics summarised
in \cref{tab:metrics}. Apart from \texttt{compute\_time\_ms} and the
positional accuracy measure \texttt{max\_waypoint\_err}, all metrics
are computed from a dense sample of 1000 evenly spaced points
generated by analytic evaluation of the interpolant.

\begin{table}[!t]
\centering
\caption{Quantitative evaluation metrics. Norms are Euclidean.}
\label{tab:metrics}
\begin{tabularx}{\linewidth}{lX}
\toprule
\textbf{Metric}     & \textbf{Definition} \\
\midrule
\texttt{max\_speed}    & $\max_t \| \dot{\vec{p}}(t) \|$ \\
\texttt{max\_accel}    & $\max_t \| \ddot{\vec{p}}(t) \|$ \\
\texttt{rms\_jerk}     & $\sqrt{\frac{1}{T} \int_0^T \| \dddot{\vec{p}}(t) \|^2 dt}$
                          (trapezoidal quadrature) \\
\texttt{path\_length}  & $\int_0^T \| \dot{\vec{p}}(t) \|\, dt$
                          (trapezoidal quadrature) \\
\texttt{max\_waypoint\_err} & $\max_i \|\vec{p}(t_i) - \vec{p}_i\|$ \\
\texttt{max\_velocity\_disc} & $\max_i \|\dot{\vec{p}}(t_i^+) - \dot{\vec{p}}(t_i^-)\|$
                                (probe with $\varepsilon = 10^{-6}(t_n - t_0)$) \\
\texttt{compute\_time\_ms}   & wall-clock time to construct and densely
                                evaluate the interpolant \\
\bottomrule
\end{tabularx}
\end{table}

The \texttt{max\_velocity\_disc} metric provides a direct,
quantitative measure of $C^1$ violation. For a $C^1$ method the
left- and right-limits of the velocity at any internal knot agree,
so the probe returns a value proportional to $2\varepsilon \cdot
\max\|\ddot{\vec{p}}\|$ (essentially numerical noise). For linear
interpolation the limits differ, and the probe returns the actual
magnitude of the velocity jump.

\subsection{Validation methodology}
\label{sec:validation-method}

To verify that the analytic-derivative formulas of each
interpolation method are correctly implemented, we densely sample
each method's position output and compute first, second, and third
derivatives by finite difference using the schemes from
Section~\ref{sec:bg-diff}. We then compare the analytic and empirical
derivatives over the interior of the trajectory (excluding the
boundary regions where finite differences fall back to lower-order
one-sided stencils). Disagreement that scales with $h^k$ for the
expected order of accuracy indicates a correctly-implemented analytic
derivative; disagreement that does not scale with $h$ would indicate
a bug.

\subsection{Boundary-condition sensitivity study}
\label{sec:bc-sensitivity-method}

After observing that the natural-boundary quintic spline produced
unexpectedly large peak acceleration and RMS jerk on the baseline
mobile-robot dataset (Section~\ref{sec:results}), we conducted an
ad-hoc sensitivity study comparing two boundary-condition variants
of the same quintic spline:

\begin{itemize}[leftmargin=*]
\item \textbf{Natural BC}: $\vec{v}_0 = \vec{a}_0 = \vec{v}_n = \vec{a}_n = \vec{0}$.
\item \textbf{Cubic-derived clamped BC}: $\vec{v}_0, \vec{v}_n, \vec{a}_0, \vec{a}_n$
      are read off the cubic-spline solution at the same knots, then
      passed to the quintic spline as clamped boundary values.
\end{itemize}

The sensitivity study quantifies how much of the observed quintic
behaviour is attributable to the boundary condition rather than the
method itself.
